\section{Round cones and the reduction of thin shapes}\label{sec:round-transfer}

The round model has two transverse directions.  Its regularity tests include
neighborhoods of null hyperquadrics.  A thin neighborhood of one such
hyperquadric becomes a cylinder after a conformal change of coordinates and a
rescaling in one transverse direction.  The purpose of this Section is to
justify this reduction, including its effect on wave packets and regularity.

\subsection{The round experiment}

\begin{definition}[Round experiment]\label{rt:model}
Write a spacetime point as $(t,b,y)\in\R\times\R^2\times\R$, and put
\[
 q(t,b,y)=ty-|b|^2,\qquad
 L_A=\partial_t+A\cdot\partial_b+|A|^2\partial_y.
\]
The classical trajectories are
$(B+tA,Y+t|A|^2)$, with arbitrary positive phase energy of total mass $e$.
The initial phase support has finite-power size, and its density relative
to Lebesgue measure is at most $e$ times a finite power of the final scale.
These bounds are uniform for the families in a fixed setup.  Passive
indices, when present, range over standard Borel spaces; separable
Hilbert multiplicity is unrestricted.
In the wave experiment the initial datum belongs to
$L^2(\R^3;\mathcal H)$, where $\mathcal H$ is any separable Hilbert space,
and has squared norm $e$.  Its evolution is
\[
 \widehat{U(t)v}(\xi,-\rho)
 =e^{-it|\xi|^2/(4\rho)}\widehat v(\xi,-\rho).
\]
The parameter $H>0$ and its reciprocal have finite-power size, and
$H\rho$ belongs to a fixed compact subinterval of $(0,\infty)$.
At length $\ell$ set
\[
 d_\ell=(H/\ell)^{1/2},\qquad q_\ell=(H\ell)^{1/2}.
\]
The analyzing operator $V_\ell$ applies, for each $A\in\R^2$, the multiplier
\[
 \chi(H\rho)\psi\bigl((\xi/(2\rho)-A)/d_\ell\bigr),
 \qquad \int_{\R^2}|\psi|^2=1,
\]
and then takes the inverse spatial Fourier transform.  Its row measure is
$d_\ell^{-2}\,dA\,db\,dy$.  Here $\psi$ is smooth and compactly supported,
and $\chi$ is smooth, supported in a positive band, and bounded by one.
A \emph{plateau analysis} has $\chi=1$ on the incoming spectrum.  Such an
analysis is isometric; every analysis is contractive.  Masks act pointwise
by contractions on the rows and are followed by synthesis $V_\ell^*$.
After synthesis, subsequent plateau cutoffs may use a fixed band enlargement.
Only finitely many cuts at fixed depths are used on any inner sequence.
Initial row energy outside a finite-power phase box is smaller than every
power of the final scale, relative to $e$.
\end{definition}

For a classical experiment put $d_\ell=0$.  In both versions, $\mu_I$ denotes
the row energy at the left endpoint of the interval $I$.  The following
normalization specifies the tests without depending on a choice of chart.
For a test at time $T$, length $\ell$, and width $r>0$, use coordinates
\begin{align*}
 u&=(t-T)/\ell,&
 x&=(b-b_*-\ell uA_*)/(\ell r),\\
 z&=(y-y_*-2A_*\cdot(b-b_*)+\ell u|A_*|^2)/(\ell r^2),&
 V&=(A-A_*)/r.
\end{align*}
The normalized rays again have velocities $(V,|V|^2)$.

\begin{definition}[Round regularity]\label{rt:regularity}
A round test imposes $|x|,|z|,|V|\le1$ at $u=0$ and has $f=1$.
A hyperquadric test additionally imposes
\[
 |F|\le f,\qquad |G|\le f,\qquad 0<f\le1,
\]
where
\[
 F=c+du+\alpha\cdot x+\beta z+k(uz-|x|^2),\qquad G=L_VF,
\]
and
\[
 |\alpha|^2-4\beta d+4kc=1,\qquad
 \max(|c|,|d|,|\alpha|,|\beta|,|k|)\le100.
\]
Every test must satisfy $rf\ge d_\ell$.  Its weight is
$w=r^4f^{1-\eta}$, where $0<\eta<1$ is fixed.
Regularity with constant $\kappa$ means $\mu_I(E)\le\kappa w(E)$ for every
test $E$.  Initial regularity is required, and $e/\kappa$ and $\kappa/e$
have finite-power size.  At a final cut one may assign measures dominated
in sum by the row energy, with each assignment supported on its own test.
Writing their masses as $m_\lambda$, the normalized functional is
\[
 \frac{s^{-1/2}}{e\sqrt\kappa}
 \sum_\lambda m_\lambda^{3/2}w_\lambda^{-1/2}.
\]
\end{definition}

Direct differentiation gives
\begin{equation}\label{rt:discriminant-identity}
 L_VG=0,\qquad L_Vg=0,\qquad
 g=\partial_VG=\alpha+2(\beta+ku)V-2kx,
 \qquad |g|^2=1+4(\beta+ku)G-4kF.
\end{equation}
Translations, boosts with their indicated vertical tilts, rotations, and
the preceding parabolic changes of scale preserve the class of experiments.
The parameter in normalized coordinates is $H/(\ell r^2)$, and the change
in regularity is exactly reciprocal to the change in weights.  Multiplying
the datum by the square root of the spatial Jacobian preserves its energy.

\begin{lemma}[Packet locality]\label{rt:packet-locality}
For $\ell'\le\ell$ and $|u|\lesssim\ell$, the kernel of
$V_{\ell'}U(u)V_\ell^*$ vanishes unless
$|A'-A|\lesssim d_{\ell'}$.  On this angular support its absolute value is
at most
\begin{equation}
 \frac{C_N}{q_\ell^2H}
 \left(1+\frac{|b'-b-uA|}{q_\ell}
 +\frac{|y'-y-u|A|^2-2A\cdot(b'-b-uA)|}{H}\right)^{-N}.
 \tag{cone-kernel}\label{eq:cone-kernel}
\end{equation}
Consequently, after any fixed finite stack, retained correspondence paths
from a row at cut $i$ to a row at cut $j$ satisfy
\begin{equation}
 \begin{split}
 |A_j-A|&\le s^{-O(\eps)}d_{\ell_j},\\
 |b_j-b^p|&\le s^{-O(\eps)}q_{\ell_i},\\
 |y_j-y^p-2A\cdot(b_j-b^p)|&\le s^{-O(\eps)}H.
 \end{split}
 \tag{acc}\label{eq:acc}
\end{equation}
Here $(b^p,y^p)$ is the prediction along the earlier row's ray.  Removing
all kernel pairs outside these correspondences changes the operators by
an arbitrarily small power error, for every prescribed $\eps>0$.
\end{lemma}

\begin{proof}
The two angular symbols have disjoint supports unless the angular assertion
holds.  On their overlap use frequency coordinates
$q_\ell(\xi-2\rho A)$ and $H\rho$.  The amplitude is compactly supported
and smooth with uniformly bounded derivatives.  The quadratic part of the
propagation phase has the same property because
$uH/q_\ell^2=O(1)$.  The remaining linear phases are precisely the two
position differences in Equation~\eqref{eq:cone-kernel}.  Integration by
parts gives that bound, including its normalization $1/(q_\ell^2H)$.
The row measures and the inverses of the packet scales have finite-power
size.  Schur's test therefore makes the omitted operator norm smaller than
any prescribed power by choosing $N$ last.

For a path through a finite sequence of cuts, an angular change at length
$\ell_m$ acts for at most $O(\ell_m)$ additional time.  It contributes
$O(q_{\ell_m})$ in transverse position and
$O(\ell_m d_{\ell_m}^2)=O(H)$ to the tilt-corrected vertical position.
Replacing an intermediate tilt by the original tilt costs at most
$O(d_{\ell_m}q_{\ell_m})=O(H)$, with the same dilation factors.  Summing
finitely many such bounds proves Equation~\eqref{eq:acc}.  Successive
truncation and contraction also show that an output region with no retained
predecessor path has negligible energy.  This last assertion uses operator
norms and does not posit a positive transition kernel for the wave.
\end{proof}

\begin{lemma}[Enlarged tests]\label{rt:covering}
Suppose $P\ge2$.  Enlarge the normalized round bounds by fixed powers of
$P$, allow coefficients bounded by fixed powers of $P$ with discriminant
one, and replace both sublevels by $P^Cf$.  The resulting set has mass at
most $\kappa r^4f^{1-\eta}P^{C'}$, provided $rf\ge d_\ell$.
The exponent $C'$ depends only on the stated enlargement exponents.
Restriction, synthesis, and plateau reanalysis at one cut consequently
preserve regularity with an arbitrarily small exponent loss.
\end{lemma}

\begin{proof}
If $f$ exceeds a sufficiently small inverse power of $P$, cover by round
tests; their number and the reciprocal of $f^{1-\eta}$ are polynomial in
$P$.  Otherwise Equation~\eqref{rt:discriminant-identity} gives
$|g|\asymp1$ on the set.  Subdivide its normalized phase box into parabolic
boxes of transverse width $a=P^{-D}$, centered on actual witnesses.
After centering at a witness, boosting to its ray, applying the parabolic
change of scale, and dividing the polynomial by $a$, its transverse
linear coefficient is $g$ at that witness.  Its vertical and quadratic
coefficients are $O(aP^C)$; its constant and central-ray derivative are
$O(P^Cf/a)$.  Choose $D$ first and the smallness threshold for $f$ second.
All coefficients then satisfy the bounds in Definition~\ref{rt:regularity},
after bounded further subdivisions.  The discriminant remains one.
The new relative thickness can be taken as $P^{C''}f/a$; thus the
absolute uncertainty width is at least $rf$.  There are polynomially
many boxes, and their total test cost has the asserted bound.

For reanalysis, Equation~\eqref{eq:cone-kernel} puts all relevant
predecessors of a test inside such an enlargement with $P=s^{-O(\eps)}$.
Apply the preceding covering to their union and use contraction on that
union.  Its complement contributes an arbitrary power error.  Finally let
$\eps$ tend to zero after the inner scale limit.
\end{proof}

All subsequent polynomial truncations are legitimate under Definition~
\ref{rt:model}.  Here is the useful quantitative reason.  For a thin test,
Equation~\eqref{rt:discriminant-identity} gives a velocity sublevel of
normalized area $O(f)$ at each allowed position.  A finite-power phase
density therefore bounds mass divided by $ew$ by
$s^{-C}r^2f^\eta$.  In the classical compact-support case, slicing instead
in the original velocity coordinates gives the additional bound
$s^{-C}r^{-3}f^\eta$ for large $r$.  In the wave case $rf\ge d_\ell$
excludes arbitrarily thin tests at a fixed large width, and the first
bound together with total mass excludes extreme widths.  Thus, above any
fixed power floor, widths, their reciprocals, and centers can be restricted
to finite-power ranges.  Summed assignment mass, rather than the number
of labels, controls the discarded small assignments.  Error estimates
after subdivision use the original energy as reference; components too
small to normalize are bounded directly by contraction and the lower
weight cutoff.

\begin{lemma}[Single-time estimate]\label{rt:single}
Relative to regularity at a parent cut, the energy in a test of widths
$(r,f)$ at any one descendant of relative length $\delta$ satisfies
\begin{equation}
 m\le s^{-o(1)}\kappa r^4 f^{1-\eta}\delta.
 \tag{single}\label{eq:single}
\end{equation}
The same conclusion holds through any fixed finite stack.
\end{lemma}

\begin{proof}
Normalize angular units by $r$, but keep the parent time and position units.
At the target time $\theta$, the test has $|b|,|y|\le\delta$, $|A|\le1$,
and, in the non-round case,
\[
 F=\alpha\cdot b+\beta y+c\delta+d(t-\theta)
       +k((t-\theta)y-|b|^2),\qquad |k|=K\le100/\delta,
\]
with $|F|\le\delta f$, $|G|\le f$.  The remaining coefficients are
bounded, and the discriminant is one.  Write $P=s^{-O(\eps)}$.
The uncertainty condition gives $d_1\le f\sqrt\delta$, and retained
predecessors satisfy
\[
 |A-A_*|\le Pf,\quad |b-b_*|\le Pd_1,\quad
 |y-y_*-2A\cdot(b-b_*)|\le Pd_1^2.
\]
Expanding first in velocity and then in position shows
$|G|\le P^Cf$ and $|F|\le P^Cf$ on the predecessor rays at the parent
time.  The potentially largest quadratic error is
$Kd_1^2\lesssim f^2$.

Cover angles by caps of width $a\gtrsim\max(d_1,\delta)$, $a\le1$.
Their number is at most
\[
 P^Ca^{-2}(K\delta+f+a).
\]
Indeed the relevant angles obey
$|d+\alpha\cdot A+\beta|A|^2|\lesssim P^C(K\delta+f)$.
When this sublevel and its $a$-enlargement are small, its gradient is
bounded below by the discriminant identity; slicing the quadratic gives
the claimed area bound.  Otherwise the full angular area suffices.
For a round test use $K=0$ and $f=1$ in these counts.

Within a cap centered at $A_0$, only $P^C$ transverse position cells of
width $a$ are needed.  The required length in $y-2A_0\cdot b$ is at most
the following quantity, with $(f+a)/K=\infty$ when $K=0$:
\[
 P^C\left(a^2+\min\{\delta,(f+a)/K\}\right).
\]
This follows by writing
$G=d+\alpha\cdot A_0+\beta|A_0|^2+k(y-2A_0\cdot b)$; the error is
$P^C(a+Kd_1a)$.  Divide the displayed length by $a^2$ to count vertical
cells.  Backflow preserves the resulting parent round boxes up to
bounded enlargement.

Put $a_0=\min(\sqrt\delta,1/K)$, with the same convention when $K=0$.
If $f\le a_0$, take $a=a_0$ and retain the hyperquadric condition.
Lemma~\ref{rt:covering}, applied after dividing the recentered polynomial
by $a$, bounds the cost of each box by
$P^C\kappa r^4a^{3+\eta}f^{1-\eta}$.
For $K\le\delta^{-1/2}$ the total cost, divided by
$\kappa r^4f^{1-\eta}$, is at most
$P^C\delta^{1+\eta/2}$; for $K>\delta^{-1/2}$ it is at most
$P^C\delta K^{-\eta}$.  Both are $O(P^C\delta)$.

If $f>a_0$, take $a=\max(a_0,d_1)$ and use only round bounds.  Their
total cost divided by $P^C\kappa r^4$ is
\[
 (K\delta+f+a)\left(a^2+\min\{\delta,(f+a)/K\}\right)
 \lesssim\delta f.
\]
For $K\le\delta^{-1/2}$ this uses $a=\sqrt\delta$ and $f>\sqrt\delta$.
For larger $K$, use $a\lesssim f$,
$a^2\lesssim K^{-2}+\delta f^2$, and $K\lesssim\delta^{-1}$.
Since $f\le f^{1-\eta}$, this proves the desired bound.
Contraction on the covered predecessors, followed by Lemma~\ref{rt:packet-locality},
handles the wave and finite-stack versions.  Letting $\eps$ decrease
gives Equation~\eqref{eq:single}.
\end{proof}

\subsection{Free gaps and regularity splitting}

For each version, let $\Gamma$ be the supremum of upper exponents of the
normalized functional for free experiments, with the closure convention
of the framework.  Widths and parameters have fixed finite-power ranges
on each inner sequence; exterior limits of their exponents are allowed.
Equation~\eqref{eq:single} and total mass over $O(s^{-1})$ cuts imply
$\Gamma\le1$, so this upper exponent is finite.

\begin{proposition}[Composition of free gaps]\label{rt:free-composition}
Inserting finitely many admissible pointwise row contractions followed
by synthesis, as in Definition~\ref{rt:model}, at fixed depths does not
increase $\Gamma$.  Repeated operations at the same cut are permitted.  Reconstruction
uses plateau analyses, while an old non-plateau terminal analysis may be
retained for coherent recombination.
\end{proposition}

\begin{proof}
At each new cut peel off successes from tests at descending dyadic
regularity thresholds $b$, from a high power multiple of the incoming
regularity down to a low power multiple.  The top threshold exceeds a
crude polynomial regularity bound.  At threshold $b$, the entering
remainder is $O(b)$-regular because the preceding threshold $2b$ has
been exhausted.  A success has mass exceeding $bw$.  All success
portions at this level are disjoint submeasures of that entering
remainder, so their union is also $O(b)$-regular, even if the peeling
returns to previously used tests.  Synthesize this union as one state
per interval and threshold, not as one state per success.  The polynomial
weight cutoff makes the process finite; the number of threshold levels
is logarithmic.  Synthesis and Lemma~\ref{rt:covering} give post-states
with regularity $s^{-o(1)}b$.  If $z$ is the raw success mass summed over
intervals, the sum of post-state energies is at most $z$.

At depth $a\in(0,1)$ the raw prefix cube sum is at least $z\sqrt b$.
The free bound on the preceding segment gives
\[
 s^{-a/2}z\sqrt b\le s^{-a\Gamma-o(1)}e\sqrt\kappa.
\]
The free bound on the remaining segment therefore bounds its final cube
sum by $s^{1/2-\Gamma-o(1)}e\sqrt\kappa$.  Summing the logarithmically
many levels costs zero exponent.  The low-floor remainder is bounded
using its tiny regularity and the total cut energy.  Applying the same
argument on each parent proves the relative statement for two interior
cuts and, by finite induction, the full assertion.

At a repeated cut, Lemma~\ref{rt:covering} supplies the necessary
regularity without a positive-length gap.  At the last cut, split before
the terminal operators and use Lemma~\ref{rt:single} together with total
mass.  For coherent recombination freeze the allocation densities
$a_\lambda\ge0$ on the output rows, with $\sum_\lambda a_\lambda\le1$
at each cut.  If $T_\lambda$ is the corresponding output map, then
\[
 v\longmapsto
 \left(\sum_\lambda w_\lambda^{-1/2}
       \|\sqrt{a_\lambda}\,T_\lambda v\|_2^3\right)^{1/3}
\]
is a seminorm by the triangle inequalities in $L^2$ and weighted
$\ell^3$.
The old terminal cutoff can consequently be retained throughout.
\end{proof}

\begin{proposition}[Saturation at inserted cuts]\label{rt:saturation}
Suppose $\Gamma>0$.  In the outer closure of homogeneous free exponent
data, let $p_0$ be the minimal mass exponent at upper exponent $\Gamma$,
and let $d$ be the associated total-mass exponent.  Then
\[
 \Gamma=d+(1-p_0)/2,\qquad p_0\ge1,\qquad 0<d\le1.
\]
Fix any finite depth list.  After realizing these data with arbitrarily
small exterior errors, every regularity splitting preserving the
terminal functional has limiting dominant regularity and summed energy
described by
\begin{equation}
 \kappa_a=\kappa s^{p_0a+o(1)},\qquad E_a=e s^{-da+o(1)}.
 \tag{CE}\label{eq:CE}
\end{equation}
In particular, an arbitrarily short fixed final gap may be taken free,
with a regularly split incoming state, when classifying terminal shapes.
Equation~\eqref{eq:CE} denotes identities of exponent data in this
outer closure; it does not assert that one inner sequence attains the
closure extrema exactly.
\end{proposition}

\begin{proof}
Homogenization gives
$m_\lambda/(\kappa w_\lambda)=s^{p_0+o(1)}$ and hence the identity for
$\Gamma$.  Equation~\eqref{eq:single} gives $p_0\ge1$, while total
energy at $O(s^{-1})$ cuts gives $d\le1$; positivity of $\Gamma$ gives
$d>0$.  The retained terminal tests have total weight at most
$s^{-(p_0+d)+o(1)}e/\kappa$.

For completeness, minimality must be applied to free gaps, not simply
asserted for masked experiments.  Start with any saturated finite stack
having terminal homogeneous exponent $p'$.  Insert ordinary regularity
splits immediately before each positive-depth group of operators,
including terminal operators when necessary.  Between successive splits,
use the post-synthesis state as input and the next raw peeling successes
as output.  This is a free experiment.  Every bound in the composition
proof must be sharp in exponent, since otherwise the terminal functional
loses a fixed power.  Thus using full raw mass instead of retained mass,
or inherited regularity instead of sharp regularity, also loses no power.
The regularity exponent increment on each free gap is at least $p_0$
times its depth length: sum the uniform free bounds over starting
intervals and select a near-maximal one, then apply the definition of
$p_0$ to its homogeneous raw successes.  The same argument applies to a
last free gap ending in the retained final assignments.  If terminal
operators occur at that depth, their output satisfies
$m_\lambda\le s^{-o(1)}b_{\rm end}w_\lambda$ by inherited regularity;
the final homogeneous exponent is therefore at least the accumulated
regularity exponent across the preceding free gaps.

To account for changes in masses under coherent splitting, denote the
new final homogeneous exponent by $p''$, and put
$d''=\Gamma+(p''-1)/2$.  Free-gap minimality gives $p''\ge p_0$.
The retained labels still have total weight at most
$s^{-(p'+d')+o(1)}e/\kappa$, where $d'=\Gamma+(p'-1)/2$.
For their new saturated masses the same total weight has exponent
$p''+d''$.  Hence $p''+d''\le p'+d'$, so $p''\le p'$.
This proves $p'\ge p_0$ without identifying the masses before and after
splitting.  Starting with $p'=p_0$ forces equality on every free gap.
The sharpness of the energy sums in Proposition~\ref{rt:free-composition}
then gives both identities in Equation~\eqref{eq:CE}.  A split originally
made without all these isolating splits satisfies the same conclusion
by applying the just-proved inequality $p'\ge p_0$ separately to its
prefix and suffix.  To implement this in the outer closure, fix the
finite depth list first and realize the extremal data to arbitrary
exterior accuracy.  Single-time bounds, total energy, and sharpness of
the composed estimates bound all added intermediate exponents.
Extract convergent subsequences of these finitely many exponents and
then let the exterior error vanish.  The preceding inequalities pass
to their limits and give Equation~\eqref{eq:CE}.  These exterior errors
are distinct from inner $o(1)$ errors.  Therefore a further fixed final
depth can be inserted as
close to the endpoint as desired, leaving its continuation free.
\end{proof}

\begin{proposition}[Thin endings reduce to cylinders]\label{rt:thin-exclusion}
Let $\Gamma_{\rm cl}^{\rm cyl}$ and $\Gamma_{\rm hyb}^{\rm cyl}$
be the classical and hybrid cylinder exponents.  A saturated round
experiment with terminal thickness of positive limiting exponent
satisfies
\[
 \Gamma\le\Gamma_{\rm cl}^{\rm cyl}
 \quad\hbox{in the classical case},\qquad
 \Gamma\le\max(\Gamma_{\rm cl}^{\rm cyl},\Gamma_{\rm hyb}^{\rm cyl})
 \quad\hbox{in the wave case}.
\]
The conclusion includes infinite limiting thickness exponents.
\end{proposition}

The proof has three reductions.  First group the thin ending tests around
one hyperquadric.  Then flatten that hyperquadric and transfer both input
regularity and terminal assignment capacities through its conformal chart.
Finally freeze the very fine base variables while retaining the normal
quantum evolution; the factors introduced by the normal rescaling will
cancel.

\subsection{Conformal charts and thin families}

We now work on a free gap of relative length $\delta=s^h$, where $h>0$
is fixed before taking the scale limit.  If the ending thickness has
strictly positive limiting exponent, $h$ can be chosen so small that
\begin{equation}\label{rt:very-thin}
 f\le\delta^Z
\end{equation}
for any prescribed fixed $Z$.  Normalize by the common ending angular
width and the parent time length.  Ending weights are then comparable
to $f^{1-\eta}$, and $d_\delta\lesssim f$.

Packet backprojections let us split into bounded parent phase charts:
use unit angular cells, tilted unit position cells, and Lemma~
\ref{rt:packet-locality}.  A label meets only boundedly many such source
cells, their raw energy budgets sum to at most $e$, and Lemma~
\ref{rt:covering} preserves input regularity.  Select a near-maximal
chart by the uniform free bound.  At positive saturation we may retain
homogeneous labels satisfying
\begin{equation}\label{rt:heavy-labels}
 m_\lambda\ge\delta^4\kappa f^{1-\eta},\qquad
 \#\{\lambda\}\lesssim
 \frac{\delta^{-5}e}{\kappa f^{1-\eta}}.
\end{equation}
Indeed $\Gamma>0$ implies that the homogeneous mass exponent of this
free gap is less than $3$, and the second bound follows by summing
cut energy.  Comparable widths produce the same conclusions.

In the bounded parent chart, write the ending hyperquadrics as
\[
 F_\lambda=c_\lambda+d_\lambda t+
 \alpha_\lambda\cdot b+\beta_\lambda y+k_\lambda q,
 \qquad n_\lambda=(c_\lambda,d_\lambda,\alpha_\lambda,
                   \beta_\lambda,k_\lambda).
\]
Their discriminant is one, $|n_\lambda|\lesssim\delta^{-1}$, and at their
readout time $\theta$ one has $|\beta_\lambda+\theta k_\lambda|\lesssim1$.
Predecessors lie in bounded sets $S_\lambda$ satisfying
$|F_\lambda(0)|,|L_AF_\lambda|\le\delta^{-1}f$, after harmless small
power enlargements.

\begin{lemma}[Flattening a null hyperquadric]\label{rt:conformal-chart}
Let $|n|\le P$, $P\ge2$, with discriminant one.  Near any bounded ray
on which $|F(0)|+|L_AF|$ is smaller than a sufficiently high inverse
power of $P$, there is a rational conformal chart carrying $F=0$ to
$b'_2=0$.  A polynomial number of such charts cover the relevant rays.
On buffered subcharts, both chart directions have derivative bounds
$C_jP^{Cj+C}$, time along a ray increases at a rate comparable to one,
and parabolic phase boxes of width $R\le1$ map into polynomial
enlargements of boxes of the same width.
The induced linear map on hyperquadric coefficients preserves the
discriminant and has norm and inverse norm bounded by powers of $P$.
\end{lemma}

\begin{proof}
Center on the given ray.  Equation~\eqref{rt:discriminant-identity} makes
its transverse coefficient nonzero.  Move its initial point in that
direction to set $F(0)=0$, and then move its velocity in that direction
to set $L_AF=0$.  The implicit scalar equations have derivatives bounded
below; the changes are polynomial in $P$ times the original errors.
Center and boost on this exact ray, then rotate so that
$F=b_2+\beta y+kq$.
First make the linear $q$-isometry
\[
 u=b_2+\beta y,\qquad t_{\rm new}=t+2\beta b_2+\beta^2y,
\]
retaining $b_1,y$.  In these coordinates put $v=(0,0,k,0)$ and
\begin{equation}
 x'=\frac{x+v q(x)}{\mathcal D(x)},\qquad
 \mathcal D(x)=1+2\langle v,x\rangle+q(v)q(x).
 \tag{conf}\label{eq:conf}
\end{equation}
Here $\langle\ ,\ \rangle$ is the polar form of $q$.
Then $b'_2=F/\mathcal D$, the inverse uses $-v$, and
$q(x')=q(x)/\mathcal D(x)$.  Differentiation gives metric scale
$\mathcal D^{-2}$.  The time axis is fixed and has no poles, so
inverse-polynomial neighborhoods of its buffered segment have all the
stated derivative and time properties.

Division by $\mathcal D$ carries a polynomial $c+lx+k'q(x)$ to
\[
 c(1-2\langle v,x'\rangle+q(v)q(x'))
       +l(x'-vq(x'))+k'q(x'),
\]
which proves linearity, invertibility, and preservation of the
discriminant.  Null lines map to null lines, as is also immediate from
the fractional-linear restriction of Equation~\eqref{eq:conf} to a
null line.  Their polar covector
$dy-2A\cdot db+|A|^2dt$ is preserved up to a controlled nonzero scalar.
On a phase box this controls the vertical displacement to first order;
Taylor's remainder is $O(P^CR^2)$.  Transverse and angular displacements
are $O(P^CR)$.  Backflow to a common time preserves these conclusions.
Subdivide the bounded phase region into sufficiently small polynomial
cells around actual witnesses to obtain the stated covering.  Boxes
larger than the chart buffer use a crude polynomial covering instead.
\end{proof}

\begin{lemma}[Overlap of thin predecessors]\label{rt:thin-overlap}
For
$K=\max(1,\min_\pm|n_\lambda\pm n_{\lambda'}|/f)$,
\begin{equation}
 \mu_0(S_\lambda\cap S_{\lambda'})
 \le\kappa\delta^{-C-o(1)}f^{1-\eta}/K.
 \tag{CI}\label{eq:CI}
\end{equation}
The fixed exponent $C$ is independent of $Z$ in Equation~
\eqref{rt:very-thin}.
\end{lemma}

\begin{proof}
Flatten $F_\lambda$ by Lemma~\ref{rt:conformal-chart}, with $P$ a fixed
power of $1/\delta$.  At time zero, $|A'_2|,|B'_2|\le\delta^{-C}f$.
Dropping these small coordinates from the second hyperquadric leaves,
in $X=A'_1,B=B'_1,Y=Y'$, the two conditions
\[
 |c'+\alpha'B+\beta'Y-k'B^2|\le\delta^{-C}f,\qquad
 |d'+\alpha'X+\beta'X^2+k'(Y-2XB)|\le\delta^{-C}f.
\]
The five remaining coefficients are at most $\delta^{-C}Kf$.
If $K>\delta^{-C}$, at least one of $\alpha',\beta',k'$ has size
at least $\delta^CKf$: otherwise nonempty intersection controls the
two constants, while discriminant one forces the omitted transverse
coefficient to be close to $1$ or $-1$, contradicting the definition
of $K$ and invertibility of the coefficient map.

Divide by $f$ and apply the cylinder covering calculation
Equation~\eqref{eq:I}.  It gives base cells of widths
$R\gtrsim1/M$, where $\delta^CK\le M\le\delta^{-C}K$, with
$\sum R^{3+\eta}\le\delta^{-C-o(1)}/K$.
For $K\le\delta^{-C}$ use a unit covering.  Since
$M\le\delta^{-C}/f$, the extra second-coordinate tilts fit inside
polynomial enlargements of full round cells.  Pull them back by the
chart.  In each cell, the first shell and Lemma~\ref{rt:covering} give
cost $\kappa\delta^{-C}R^{3+\eta}f^{1-\eta}$, using relative thickness
$f/R$.  If $R<f$, the ratio $f/R$ is polynomially bounded and round
tests give the same conclusion with polynomial loss.  Their uncertainty
widths may be enlarged by that same factor.  Summing proves the result.
\end{proof}

For a fixed label, partners at dyadic coefficient distance $K$ number
at most $\delta^{-C}K^{1-\eta}$.  Their predecessor union lies in the
first shell enlarged to thickness $\delta^{-C}Kf$.  Regularity and
contraction bound assignable mass at each time by
$\kappa\delta^{-C}(Kf)^{1-\eta}$; sum over times and use
Equation~\eqref{rt:heavy-labels}.  Equation~\eqref{eq:CI} and the total
label count now give
\[
 \mu_0\{\text{a pair with }K>\delta^{-D}\}
 \le e\delta^{-C}\sum_{K>\delta^{-D}}K^{-\eta}.
\]
Choose $D$ so that this is a positive power smaller than $e$.
Restriction and synthesis preserve regularity; the uniform free bound
makes the removed state's final contribution negligible at saturation.

Choose $D_1>D+5$ and put
\begin{equation}\label{rt:group-thickness}
 f_* =\delta^{-D_1}f.
\end{equation}
Round one attached coefficient vector, up to sign, on an $f_*$ lattice
to partition the remaining rows into groups.  Each label connects to
boundedly many groups.  Choose an actual discriminant-one representative
$F_L$ in each group.  All connected coefficients differ from it by
$O(f_*)$, and all initial rows satisfy
$|F_L(0)|+|L_AF_L|\lesssim f_*$.
With raw budgets $e_L$ one has
\begin{equation}\label{rt:group-budgets}
 \sum_L e_L\le e,\qquad
 e_L\le\kappa\delta^{-C}f_*^{1-\eta}.
\end{equation}

Large quadratic coefficients cannot carry saturation.  Indeed, if
$|k_L|>\delta^{-b}$, connected times satisfy
$|\beta_L+k_L\theta|\lesssim1$ and lie in boundedly many intervals of
length $\delta^{b/2}$.  Split at that length.  On the prefix the
single-time estimate and the bounded number of active intervals give
exponent zero, while the suffix has exponent at most $\Gamma$.
Composition therefore gains $\delta^{b\Gamma/2}$ against a saturated
gap.  Sum this bound using Equation~\eqref{rt:group-budgets} before
subdividing the other groups.  We may thus retain only
$|k_L|\le\delta^{-b}$; all other coefficients are then $O(1+|k_L|)$,
by centering at a connected label and using its coefficient bounds.
Lemma~\ref{rt:conformal-chart} partitions each retained group into at
most $\delta^{-Cb}$ buffered charts.  Their norm and overlap costs are
arbitrarily small powers when $b$ is sufficiently small.

\subsection{Transporting a thin wave through a chart}

Fix one low-coefficient buffered chart, with raw budget still denoted
$e_L$, and write $P=\delta^{-Cb}$.  Normalize the transformed time range
to $[-1,3]$ and round transformed terminal event times to the
$\delta$-lattice.  At the ideal cylinder input, physical packet length
$4$ gives its normalized root analyzer; replacing $1$ by $4$ below only
changes fixed constants.  Shrink the phase patch by polynomial factors
of $P$ so that throughout a slightly longer interval all relevant rays
have valid inverse charts, positive time derivative comparable to one,
and inverse-polynomial spatial buffers.  Pulled incident covectors then
have $\rho'>0$, $\rho'\asymp\rho$, and transformed angle in a small cone
about the time axis.  Choose a larger safe cone with another polynomial
buffer.  These choices are possible by Lemma~\ref{rt:conformal-chart};
all order-$j$ cutoff derivatives are bounded by $C_jP^{Cj+C}$.

The algebraic covariance alone does not identify the transformed
$L^2$-energy.  We obtain that comparison from packet kernels.

\begin{lemma}[Conformal wave transport]\label{rt:wave-transport}
Let $u$ be the free cone wave synthesized from the retained patch.
There is an exact free cone wave $u'$ in the transformed variables,
with a fixed positive radial band. At packet length $\ell=1$ or $\delta$,
write $p\sim_\ell p'$ when the old row $p$ and transformed row $p'$
are close to the mapped and backflowed ray in angle, transverse position,
and tilt-corrected vertical position, with respective widths
$\delta^{-\gamma}d_\ell$, $\delta^{-\gamma}q_\ell$, and
$\delta^{-\gamma}H$; for terminal rows, include the corresponding chart
event time rounded to the $\delta$-lattice. Put
$N_\ell(E')=\{p:p\sim_\ell p'\text{ for some }p'\in E'\}$.
Let $\mu_{\rm in}$ be the retained raw old row-energy measure at the
parent input, and let $\mu'_{\rm in}$ be the transformed analyzed
row-energy measure at time $-1$. For every prescribed $\gamma>0$ and
$N>0$, if $b$ and the auxiliary localization losses are sufficiently
small and $Z$ sufficiently large, the following hold:
\begin{enumerate}[label=\textup{(\roman*)}]
\item for every Borel set $E'$ of transformed input rows,
\[
 \mu'_{\rm in}(E')\le
 \delta^{-\gamma}\mu_{\rm in}(N_1(E'))+C_N\delta^N e_L;
\]
\item for any finite old terminal fractional assignments of masses
$m_\lambda$, there are fractional assignments to $u'$ at the rounded
times, each supported on rows associated with its old label, whose
fractions sum to at most one on every transformed row. If $m'_\lambda$
is the mass aggregated over the rounded times for label $\lambda$,
then
\[
 m_\lambda\le\delta^{-\gamma}m'_\lambda+\epsilon_\lambda,
 \qquad \sum_\lambda\epsilon_\lambda\le C_N\delta^N e_L.
\]
\end{enumerate}
The predecessor sets are measured in the completion of $\mu_{\rm in}$.
\end{lemma}

\begin{proof}
\emph{Equation and spectral localization.}
The classical conformal covariance of the four-dimensional wave equation
\cite[Section~2]{Bateman1909} takes the following form in our coordinates.
Write $\Box=4\partial_t\partial_y-\Delta_b$.
For Equation~\eqref{eq:conf}, put $w=\mathcal D^{-1}$.  Direct
differentiation gives
\[
 \Box\mathcal D=8q(v),\qquad
 q^{-1}(d\mathcal D,d\mathcal D)=4q(v)\mathcal D,
 \qquad \Box w=0.
\]
Similarly $\Box(wx')=0$, using
\[
 q^{-1}(d\mathcal D,d(x+vq(x)))
       =2v\mathcal D+2q(v)(x+vq(x)).
\]
The differential metric identity therefore gives
\begin{equation}\label{rt:conformal-covariance}
 \Box(wF\circ\Phi)=w^3(\Box F)\circ\Phi.
\end{equation}
Apply this to the inverse chart and multiply by a spatial cutoff equal
to one on a buffered neighborhood of all transformed rays.  Denote the
result by $\Psi$.

Off the original packet flow by an inverse-polynomial buffer, every
fixed derivative of $u$ is $O_N(H^N\sqrt{e_L})$, for arbitrary $N$,
with integrable weights at infinity.  To verify this, express $u$ as
initial packet synthesis and use Equation~\eqref{eq:cone-kernel} at
length one.  Cauchy--Schwarz in the bounded initial row region costs
only a fixed power of $H^{-1}$; a fixed derivative costs another fixed
power.  The packet width $d_1=\sqrt H$ is smaller than the buffer by a
power when $Z$ is large.  Arbitrary integration order then absorbs
these costs.  Thus cutoff derivatives occur only where $u$ has this
decay, and Equation~\eqref{rt:conformal-covariance} proves the same
error estimate for $\Box\Psi$ in every fixed spatial Sobolev norm.

Choose a smooth safe-cone and radial-band projector $\mathcal P$ that
equals one on all incident pulled covectors throughout the time interval,
and evolve $\mathcal P\Psi(-1)$ exactly to define $u'$.
Uniformly for every $t'\in[-1,3]$,
$\|(1-\mathcal P)\Psi(t')\|_{H^j}=O_N(H^N\sqrt{e_L})$ for every
fixed $j,N$.
Indeed insert the old compact-band Fourier representation of $u$.
Outside the plateau the spatial phase gradient is separated from
the output frequency $\zeta'$ by
$cP^{-C}(H^{-1}+|\zeta'|)$, while its order-$j$ derivatives are at most
$C_jP^{Cj+C}H^{-1}$.  Repeated integration by parts gives arbitrary
powers of $H$, including integrable decay for $|\zeta'|\gg H^{-1}$.
The exponents of $P$ are linear in the total differentiation order,
so the required choice of $Z$ is independent of that order.
On the retained band the forced equation is first order in $t'$ after
division by $-4i\rho'$; unitarity and integration in time show
\[
 \sup_{-1\le t'\le3}\|u'(t')-\Psi(t')\|_{H^j}
       =O_N(H^N\sqrt{e_L})
\]
for every fixed $j,N$.  Sobolev embedding also gives the pointwise
comparison needed below.  Smooth band projection and bounded-time
propagation leave negligible analyzed energy outside a bounded spatial
region and the enlarged safe angular cone, by the same phase argument.

\smallskip\noindent
\emph{Packet association.} Put $Q_\ell=q_\ell^2H$.
At input, substitute the actual old synthesis into the weighted
pullback and then analyze $u'(-1)$.  At output, synthesize $u'(T')$,
propagate to the image event, pull back, and analyze on the old slice.
The preceding approximation justifies both substitutions with arbitrary
power errors.  Each resulting kernel has absolute value at most
$P^C/Q_\ell$: the analyzing kernel has size $Q_\ell^{-1}$ and
integrable Schwartz profile in tilted position units $(q_\ell,H)$,
and a propagated source atom is bounded by $C/Q_\ell$ from its
Fourier support.

Choose small $\eps_2\gg\eps_1\gg b$ and first restrict integration
to the target packet core enlarged by $\delta^{-\eps_1}$.
Its complement has arbitrary decay by the Schwartz bound.  If the
source angle differs from the image target angle by more than
$\delta^{-\eps_2}d_\ell$, express the source atom by plane waves.
In a transverse variable $(b-b_p)/q_\ell$, holding
$y-2A_p\cdot b$ fixed, the phase derivative has magnitude at least
$cP^{-C}\delta^{-\eps_2}$.  Its variation over the retained core is
at most $P^C\delta^{-C\eps_1}$, because $q_\ell^2/H=\ell\le1$.
An order-$j$ higher derivative is bounded by
$C_jP^{Cj+C}q_\ell^j/H$.  Apply the integration-by-parts operator
$(i\partial_v\phi)^{-1}\partial_v$ in a direction nonstationary at
the center.  The product rule bounds the resulting terms by
\[
 \frac{P^{C_0}}{Q_\ell}
 \delta^{N(\eps_2-C\eps_1-Cb)-C_0}
\]
times finite-power normalization factors, for arbitrary $N$.
Take $\eps_2>C\eps_1+Cb$.  The inner choice of $N$ then absorbs all
fixed polynomial row-volume costs, even when the exponent of $H^{-1}$
is very large.  This proves angular localization without any loss
proportional to that exponent.

Once angles match, use the propagated source packet envelope for
spatial localization.  The relevant propagation duration is at most
$\ell\delta^{-O(\eps_1+b)}$.  The chart sends the target core into
the associated source core: the polar covector cancels first-order
vertical variation and the remaining quadratic error is controlled by
$q_\ell^2\le H$.  Position-induced angle variation is controlled by
$q_\ell\le d_\ell$.  Backflow through the allowed time discrepancy
adds transverse error $O(\ell d_\ell)=O(q_\ell)$ and corrected
vertical error $O(\ell d_\ell^2)=O(H)$, with the same small-power
enlargements.  Applying the inverse chart proves the column version
of this assertion, including after event-time rounding.

\smallskip\noindent
\emph{Transfer of assignments.}
For fixed source and target time indices, row and column association
volumes in normalized row measure are therefore at most
$Q_\ell\delta^{-O(\eps_2)}$.  Schur's test with the kernel bound proves
the input energy comparison.  For output assignments, let
$h_\lambda(p)$ be their old fractions, including the time partition.
Pointwise Cauchy--Schwarz bounds an old assigned mass by a small-power
multiple of transformed energy with fractions
\[
 Q_\ell^{-1}\int_{p\sim p'}h_\lambda(p)
                  d_\ell^{-2}\,dA_p\,db_p\,dy_p.
\]
For a fixed $p',T'$ these fractions sum to at most a small inverse
power.  The column volume gives this for each old time.  To count those
times, inverse-map the source center event at $T'$ and call its old
time $t_0$.  A mapped associated old event differs from that center
by $O(\delta)$ in time, from rounding, and by at most
$\delta^{-O(\eps_2)}(q_\delta+H+\delta)$ in ordinary spatial
coordinates; the last term accounts for bounded velocity during
rounding.  Since $q_\delta=\delta d_\delta\lesssim\delta$ and
$H=\delta d_\delta^2\lesssim\delta$, the inverse chart gives
$|t_{\rm old}-t_0|\le\delta^{1-O(\eps_2)}$.
Old cuts have spacing $\delta$, so only
$\delta^{-O(\eps_2)}$ indices occur.  This bound is independent of
which label assigned the event to that rounded time.
Divide all fractions by the resulting bound to legitimate the new
assignments.  A fixed old label uses only polynomially many in $P$
rounded times, because its position cell has diameter $O(\delta)$;
H\"older's inequality bounds this aggregation cost by another small
power.  Choosing the losses so their combined exponent is below
$\gamma$ proves both conclusions.
\end{proof}

\begin{lemma}[Allocation from cell capacities]\label{rt:capacity}
Let $x_i$ and $y_j$ be energies in two finite measurable cell partitions,
and let $A,D\ge1$ and $\epsilon_i\ge0$.
Suppose
\[
 x_i\le A\sum_{j\in\mathcal N(i)}y_j+\epsilon_i,
\]
and every source cell $j$ belongs to at most $D$ sets $\mathcal N(i)$.
Assignments $m_{\lambda i}$ with $\sum_\lambda m_{\lambda i}\le x_i$
can, after removing total mass at most $\sum_i\epsilon_i$, be
transferred to the $y$-cells with every retained label losing at most
a factor $AD$.  Contributions from different cells may be aggregated
back into the same label without an additional factor.
\end{lemma}

\begin{proof}
Reduce the assignments proportionally in each target cell until their
sum is at most $A Y_i$, where $Y_i=\sum_{j\in\mathcal N(i)}y_j$.
The removed mass is at most $\epsilon_i$.
For $Y_i>0$, allocate
$m_{\lambda i}y_j/(ADY_i)$ to cell $j\in\mathcal N(i)$.
Its total over labels and incident $i$ is at most $y_j$.
Summing over $j$ gives exactly $m_{\lambda i}/(AD)$, so aggregation
over $i$ preserves the stated label mass.  Cells with $Y_i=0$ have
no retained assignments.
\end{proof}

\subsection{Regularity after flattening}

In transformed coordinates use
\[
 X=A'_1,\quad B=b'_1,\quad Y=y',\quad
 N=A'_2/f_*,\quad D=b'_2/f_* ,\qquad H_{\rm flat}=H/f_*^2.
\]
Exact ray velocities in $(B,Y,D)$ are
$(X,X^2+f_*^2N^2,N)$.  Initial and final normal coordinates satisfy
$|D|,|N|\le\delta^{-C\gamma}$, modulo dispensable tails; this follows
from $F_L,G_L$, conformal flattening, and Lemma~\ref{rt:wave-transport}.
The ideal cylinder replaces $X^2+f_*^2N^2$ by $X^2$.

For waves, let $(\xi,-\rho,p_N)$ be the Fourier variables dual to
$(b'_1,y',D)$, and put $\sigma=2H\rho$.  Evolution over elapsed time
$t_0$ factors exactly as
\begin{equation}\label{rt:flat-propagator}
 U_b=e^{-it_0\xi^2/(4\rho)},\qquad
 U_\sigma=e^{-it_0H_{\rm flat}p_N^2/(2\sigma)}.
\end{equation}
Thus the normal evolution has parameter $H_{\rm flat}$ even when
$H$ is much smaller than the base observation scales.  The comparison
below replaces the base evolution by exact cylinder indices while
retaining this normal propagator whenever its quantum scale is not
negligible.

\begin{lemma}[Input and terminal weights]\label{rt:weight-comparison}
Fix a finite precision $L_0$.  Consider ideal input cylinder tests with
base and normal widths $R,F_0\in[\delta^{L_0},1]$, and shear
coefficients bounded by $\delta^{-L_0}$.  Suppose
$f_*F_0\gtrsim d_1$.  Their associated old input rows have mass at most
\begin{equation}
 \delta^{-C\gamma}\kappa f_*^{1-\eta}R^{3+\eta}F_0^{1-\eta}.
 \tag{lift}\label{eq:lift}
\end{equation}
At output, split each connected old assignment by chart and rounded
time.  Each resulting assignment fits an ideal cylinder test of weight
at most
\begin{equation}
 \delta^{-C\gamma}(f/f_*)^{1-\eta}.
 \tag{endlift}\label{eq:endlift}
\end{equation}
The assertions remain true under sufficiently small fixed-power base
errors and normal errors of size
$\delta^{-C\gamma}\max(\delta^S,d_\ell/f_*)$ in velocity and
$\ell$ times this in position, where $\ell=1$ or $\delta$ and
$S$ is sufficiently large.  The constants are independent of $Z$
once $Z$ is sufficiently large relative to these fixed precisions.
\end{lemma}

\begin{proof}
For an input cylinder test, subtract $f_*$ times its position-shear
polynomial from $b'_2$, replacing the base quadratic by
$q'=t'y'-|b'|^2$.  This produces a full hyperquadric.  On matched rays,
its position and null-derivative sublevels have size at most
$\delta^{-C\gamma}f_*F_0$.  Terms introduced by the second-coordinate
square and by $f_*^2N^2$ are bounded by a fixed power times $f_*^2$;
Equation~\eqref{rt:very-thin}, with $Z$ sufficiently large, makes
them smaller than the stated sublevels.  Its discriminant is close
to one, so normalize it to one.  The full round box has width $R$:
the absolute second-coordinate widths are much smaller than $R$.
Pull back and backflow to old time zero.  Lemma~\ref{rt:conformal-chart}
costs only powers of $P$.  Center at a witness, shrink by $R$, and
divide the polynomial by $R$.  Its coefficients are bounded by small
inverse powers; the constant and time coefficients use the two
sublevels.  Lemma~\ref{rt:covering} with relative thickness
$f_*F_0/R$ now gives
\[
 \delta^{-C\gamma}\kappa R^4(f_*F_0/R)^{1-\eta},
\]
which is Equation~\eqref{eq:lift}.  The permitted errors and packet
uncertainty are absorbed by its displayed enlargement.

At output, each connected polynomial differs after flattening from
$b'_2$ by coefficients $O(P^Cf_*)$.  At its image event its position
sublevel is $O(P^C\delta f)$ and its derivative sublevel is
$O(P^Cf)$.  Propagation to the rounded time preserves these bounds.
Normalize the $b'_2$ coefficient to one, divide by $f_*$, and drop
the $O(f_*^2)$ terms just discussed.  The resulting cylinder normal
width is $\delta^{-C\gamma}f/f_*$, with shear coefficients bounded
by $\delta^{-C\gamma}$.  Its base width is $\delta^{-C\gamma}$:
the old position cell and time rounding give positions within
$\delta^{1-C\gamma}$, and angles stay bounded.  This proves
Equation~\eqref{eq:endlift}.  These tests depend only on the old label,
chart, and rounded time, so the cell allocation below can aggregate
back into their original labels.  All allowed comparison errors fit
their enlarged bounds.  Since $d_\delta\lesssim f$, the hybrid
uncertainty condition also holds.
\end{proof}

\smallskip\noindent
\emph{Completion of input regularity.}
In each ideal input constructed below, include an average over
independent exact base and normal translations and boosts of size
$\delta^l$, with $l$ fixed larger than the terminal precision
requirements.  Retain the copies as passive indices and enlarge
terminal tests.  The exact cylinder symmetries preserve the shear
family.  This smoothing extends the moderate-test comparison to full
regularity: uniformly in shear parameters,
\[
 \mu(\text{test})\le
 \delta^{-C(l)}e_L\min(R^4,1)\min(F_0^2,1).
\]
In the hybrid case the normal smoothing uses Weyl translations and the
same estimate holds conditionally at each base point.  Together with
Equation~\eqref{rt:group-budgets}, this proves the required regularity
outside a finite-power width range, by taking its cutoff exponent
large depending on $l,\eta$.  Bounded base support handles $R>1$;
used normal coordinates can be covered by unit zero-shear boxes for
$F_0>1$, at small-power cost.  Polynomial sublevel slicing against the
jitter density discards tests whose coefficients exceed a sufficiently
large finite power.  The cylinder truncation argument therefore reduces
full regularity to Equation~\eqref{eq:lift}, with $L_0$ chosen after
these cutoffs.  Compact base support and negligible normal tails are
retained in the constructions below.

\smallskip\noindent
\emph{Additive errors.}
All cell comparisons allow a prescribed error $\delta^Je_L$.
Equation~\eqref{rt:group-budgets} and the finite-power relative weights
allow $J$ to be fixed so that these errors are negligible in the final
cube sum.  When several labels share a cell, remove its excess once,
proportionally over labels, as in Lemma~\ref{rt:capacity}; do not sum an
error separately over labels.  Components of extremely small relative
energy use the direct contraction bound instead of a normalized
exponent estimate.

\subsection{An almost classical base}

Classically, use the transformed rays at time $-1$ as ideal cylinder
input, then evolve with the ideal velocities.  The omitted vertical
displacement is $O(f_*^2\delta^{-C\gamma})$, smaller than every fixed
comparison precision once $Z$ is large.  Lemma~\ref{rt:weight-comparison},
the jitter, and the energy budget therefore give the desired cylinder
comparison directly.

For waves, first fix a large $M$ depending on all comparison precisions.
If $H_{\rm flat}<\delta^M$, use the transformed initial analyzed energy
as a positive measure on ideal rays.  Partition phase into cells of a
fixed power size, smaller than the comparison tolerances but larger
than every rescaled packet width.  Equation~\eqref{eq:cone-kernel}
bounds each backflowed output cell by neighboring input capacities,
with bounded overlap and negligible error.  Lemma~\ref{rt:capacity}
transfers assignments while retaining their labels.  Equation~
\eqref{eq:lift} applies because $d_1/f_*\ll\delta^{L_0}$.
Thus this regime reduces to the classical cylinder as well.

It remains to consider
\begin{equation}\label{rt:normal-quantum-range}
 \delta^M\le H_{\rm flat}\lesssim\delta(f/f_*)^2.
\end{equation}
The base parameter $H=f_*^2H_{\rm flat}$ is still arbitrarily small
relative to every fixed comparison precision.

\begin{lemma}[Classicalizing the base]\label{rt:classical-base}
In the range Equation~\eqref{rt:normal-quantum-range}, the transformed
wave can be replaced, for the input and terminal comparisons of
Lemma~\ref{rt:weight-comparison}, by a fixed finite direct sum of ideal
hybrid cylinder states.  Their summed energy is at most
$\delta^{-C\gamma}e_L$, each regularity is at most
$\delta^{-C\gamma}\kappa f_*^{1-\eta}$, and their assignments
satisfy Equation~\eqref{eq:endlift}, with arbitrarily small total
additive errors.  The exponent $C\gamma$ is independent of the
finite-power exponent of $H^{-1}$.
\end{lemma}

\begin{proof}
\emph{Construction of the hybrid states.}
Use $L^2$-normalized coordinates $(b'_1,y',D)$ for
$g=u'(-1)$, with Fourier variables and propagators as in
Equation~\eqref{rt:flat-propagator}.  Put $X=\xi/(2\rho)$ for the
base Fourier velocity.
Smoothly restrict to $|H_{\rm flat}p_N|\le\delta^{-C\gamma}$,
with slack.  The removed norm is arbitrarily small relative to
$\sqrt{e_L}$: integrate the initial analyzed angular tail and apply
Plancherel.  The spatial normal tail is likewise negligible by packet
synthesis; the smooth frequency cutoff preserves it with slack.
Let $\Pi$ be a further real cutoff equal to one on the retained support,
so $\Pi g=g$.  All base positions are bounded up to negligible tails.

Choose observation cells with base sides $\delta^{S_1}$ and normal
velocity and position sides
\[
 \nu_\ell=\max(\delta^{S_1},\sqrt{H_{\rm flat}/\ell}),
 \qquad \ell\nu_\ell,\qquad \ell=1,\delta.
\]
Here $S_1$ exceeds the required comparison precision.  There are
polynomially many cells.  Allow a per-cell error $\delta^{J_1}e_L$,
with $J_1$ large compared with their count exponent.

Choose $K\gg M+S_1+J_1$.  Split the base first by smooth compact
square partitions in $(X,\sigma)$ and then by smooth compact square
partitions in $(b'_1,y')$, both of side $\delta^K$.
The spectral partition is a tensor product with a common
$\sigma$-lattice.  Write
\[
 W_j=M_j^{\rm pos}P_j^{\rm freq},\qquad
 (X_j,\sigma_j,B_j,Y_j)
\]
for its operators and cell representatives.  Spectral partition kernels
have spatial width $O(H\delta^{-CK})$, so restriction to bounded base
centers loses negligible energy when $Z$ is large.  Before that
restriction, $\sum_jW_j^*W_j=\Id$ exactly; afterwards the identity holds
on $g$ up to negligible error.  The column map $W$ is contractive.
Give fiber $j$ the normal input $W_jg$, taking its base spatial variable
as an auxiliary Hilbert index.  Evolve its representative classically,
and evolve its normal datum with $U_{\sigma_j}$.  This defines one
direct-sum state for both input and output measurements.  Base partitions
commute with normal cutoffs, so its normal tails satisfy the required
bounds; Gaussian packet tails allow the stated slack.

To meet the cylinder convention $1\le\sigma\le2$, partition these
passive fibers into finitely many fixed bands $a\le\sigma_j\le2a$.
For every fiber in that band use, from the outset, cylinder parameters
$(\sigma_j/a,H_{\rm flat}/a)$ and its corresponding Gaussian analyzer.
The propagator is unchanged because $H_{\rm flat}/\sigma_j$ is
unchanged.  The Gaussian position width is now
$\sqrt{H_{\rm flat}\ell/a}$; its angular width also differs from the
displayed comparison scales only by a fixed factor.  Thus all kernel
and cell estimates below hold with these analyzers, with constants
depending on the fixed bands.  At the final cylinder estimate, energy
budgets add, each band inherits the same regularity bound, and
aggregation over the finitely many bands costs a fixed factor.
The band count and the factors $a$ are fixed on the inner setup, so
there is no scale-exponent loss.

\smallskip\noindent
\emph{Cross-analysis and radial freezing.}
Let the exact and ideal analyzing maps at a measured time be
\[
 E=V^{\rm cone}_\ell U_bU_\sigma,\qquad
 I=\{V^{j,\rm flat}_\ell U_{\sigma_j}W_j\}_j.
\]
Use a plateau cone analyzer, so $E^*E=\Id$ on the retained band.
Flat analysis is isometric, giving $I^*I=\sum W_j^*W_j$.
For $C_0=E\Pi I^*$ we therefore have
\[
 C_0Ig=Eg+O(\delta^{J_2}\sqrt{e_L}),\qquad C_0^*Eg=Ig,
\]
where $J_2$ can be prescribed before the localization order is chosen.
The reference energy is the raw budget $e_L$; no lower bound for
$\|g\|^2/e_L$ is asserted.  All three operators are contractive.
Moreover $C_0$ differs in norm by at most
$C\delta^{K-M-C\gamma}$ from
\begin{equation}
 V_\ell^{\rm cone}U_b\Pi
       \sum_jW_j^*(V_\ell^{j,\rm flat})^*.
 \tag{cap-cross}\label{eq:cap-cross}
\end{equation}
To prove this, commute $U_{\sigma_j}$ through $W_j^*$, which operates
only in the base variables.  On the outer spectral support of $W_j^*$,
$|\sigma-\sigma_j|\lesssim\delta^K$.  On $\Pi$,
$H_{\rm flat}p_N^2\lesssim\delta^{-M-C\gamma}$, so
\[
 |U_\sigma U_{\sigma_j}^*-1|
       \lesssim\delta^{K-M-C\gamma}.
\]
Group the indices with common $\sigma_j$ before summing: within each
group use contraction of $W^*$; between groups use bounded spectral
overlap of their outer projections.  This proves the norm estimate
without a factor from the number of spatial or angular bins.

\smallskip\noindent
\emph{Off-diagonal bounds.}
We prove that Equation~\eqref{eq:cap-cross} has negligible norm
between an observation cell and the complement of its enlarged
neighbors.  Compress to bounded base coordinates and
$\delta^{-C\gamma}$-bounded normal coordinates; the removed energy of
$Eg$ and $Ig$ is negligible by the established localization and flat
packet propagation.  Angular mismatch in $X$ is excluded by the outer
spectral bin and cone symbol.  The latter selects
$H_{\rm flat}p_N/\sigma$ within width
$\sqrt{H_{\rm flat}/\ell}$, whereas the flat synthesis atom has a
Gaussian Fourier profile centered at
$H_{\rm flat}p_N/\sigma_j=N_j^{\rm row}$ with the same width up to
fixed factors.  The difference of $\sigma$ and $\sigma_j$ is much
smaller than observation precision.  Normal velocity mismatch outside
enlarged neighbors thus has arbitrary Gaussian gains.

For position comparison, the inverse Fourier kernel from auxiliary
base point $(b_0,y_0)$ and normal center $D_0$ has phase
\[
 \xi(b'_1-b_0)-\rho(y'-y_0)
       -t_0\xi^2/(4\rho)+p_N(D-D_0).
\]
The point $(b_0,y_0)$ is in its spatial bin.  A base discrepancy from
$(B_j+t_0X_j,Y_j+t_0X_j^2)$ beyond the neighboring observation cells
therefore gives a phase gradient of size at least $c\delta^{S_1}$
in a base frequency direction.  In coordinates $(H\xi,H\rho)$ the
large parameter is $H^{-1}$.  Higher derivatives of the factored
phase are bounded, and an order-$n$ symbol derivative costs at most
\[
 C_n\max(\delta^{-K},\delta^{-C\gamma}d_\ell^{-1})^n.
\]
This includes differentiation through the normal angular dependence
on $\rho$.  Differentiating the reciprocal phase gradient also costs
powers of $\delta^{-S_1}$; these are absorbed by the
$\delta^{-K}$ term because $K>S_1$.  Each integration by parts
therefore gains a power of
\[
 H\delta^{-S_1}\max(\delta^{-K},\delta^{-C\gamma}d_\ell^{-1}),
\]
which is a positive power of $\delta$ when $Z$ is sufficiently large,
since $Hd_\ell^{-1}\le\sqrt H$.  Its positive exponent can be fixed
before the integration order.  Arbitrary iteration absorbs the
finite-power normalizations and integration volumes.

For normal-position mismatch use frequency $H_{\rm flat}p_N$.
Symbol and Gaussian derivatives cost at most
$C_n(\ell/H_{\rm flat})^{n/2}$, with integrable Gaussian remainders.
The linear phase derivative has size
$|D-D_0|/H_{\rm flat}$, greater by a small inverse power outside the
enlarged normal cell.  Repeated integration by parts supplies arbitrary
powers of that enlargement.  Gaussian tails do the same for normal
velocity mismatch.  Crude integration and Schur's test convert these
pointwise bounds into the claimed off-diagonal norm estimate: the row
measures, auxiliary spatial supports, packet normalizations, and bin
counts all have fixed finite-power size on the inner sequence.  Choose
the order after these powers.  The geometric smallness conditions are
unchanged, because derivative exponents grow at most linearly in order.

\smallskip\noindent
\emph{Capacity and regularity comparisons.}
Contraction of $C_0$ and its reverse, the freezing error, and
the off-diagonal bounds give each cell energy at most a bounded
multiple of the opposite neighboring capacities plus
$\delta^{J_1}e_L$.  A cell belongs to only
$\delta^{-C\gamma}$ enlarged neighbor sets.  Apply
Lemma~\ref{rt:capacity} and aggregate into the original terminal labels;
there is no loss from the total number of observation cells.
Their neighbors fit the errors in Lemma~\ref{rt:weight-comparison}.
At input the reverse comparison, summed over moderate test cells,
proves Equation~\eqref{eq:lift}.  Choose $J_1$ above the cell count
and $K$ above the required freezing precision.  The jitter extension
then proves full cylinder regularity.  This gives the asserted energy,
regularity, and terminal assignments for the same finite direct-sum
state at both measurements.
\end{proof}

\begin{proof}[Proof of Proposition~\ref{rt:thin-exclusion}]
Suppose there is a strict gap from the indicated cylinder exponents.
Proposition~\ref{rt:saturation} supplies a short free final gap with
regular incoming states.  Choose its fixed depth length so that
Equation~\eqref{rt:very-thin} holds at every required precision.
The preceding reductions discard contributions that cannot saturate,
then decompose into low-coefficient buffered patches with summed raw
energy at most a small inverse power times $e$.

For each patch, Lemmas~\ref{rt:wave-transport},
\ref{rt:weight-comparison}, and \ref{rt:classical-base}, or the
classical construction, give an ideal cylinder input of energy at most
$\delta^{-C\gamma}e_L$, regularity at most
$\delta^{-C\gamma}\kappa f_*^{1-\eta}$, and terminal tests of weight
at most $\delta^{-C\gamma}(f/f_*)^{1-\eta}$.
Transferred masses lose at most a small power.  The cylinder estimate
thus gives, for either relevant cylinder exponent $\Gamma_c$,
\[
 \sum_\lambda \frac{m_\lambda^{3/2}}{f^{(1-\eta)/2}}
 \le\delta^{1/2-\Gamma_c-C\gamma-o(1)}e_L\sqrt\kappa.
\]
The factor $f_*^{(1-\eta)/2}$ from input regularity cancels its
reciprocal from the terminal weights.  Sum over patches and use the
raw energy budgets and bounded label connections.  Choose all
small-power losses below the strict exponent gap.  The resulting
bound contradicts saturation of the free round gap.

Here is the order of choices, needed to ensure uniformity.  Fix the
group threshold $D_1$, the jitter precision, and the moderate-test
cutoffs.  Next fix observation-cell and additive-error precisions,
the normal threshold $M$, and the base partition precision $K$.
Choose geometric and localization loss exponents small relative to
the strict gap, then take $Z$ large enough for every fixed precision.
Finally choose the positive final depth length $h$ sufficiently small.
After fixing the inner sequence and its finite-power bounds, choose
the integration orders; these orders remain fixed as the inner scale
tends to zero.  Their derivative
costs have linear exponents, as proved above, so they do not alter
the earlier small-loss choices.  A positive infinite thickness
exponent also satisfies every prescribed finite $Z$, proving that
case as well.
\end{proof}
